\chapter{Glossary of Symbols and Terms}
\label{ch:glossary}

This glossary collects the notation and vocabulary the book uses. The first section
lists every mathematical symbol with its meaning and the chapter that introduces it;
the second defines the architectural terms in plain language. Nothing here is new: it
is a single place to look up a symbol or a name encountered anywhere in the preceding
chapters.

\section{Symbols}
\label{sec:glossary-symbols}

\begin{table}[htbp]
\centering\small
\begin{tabular}{@{}p{2.4cm}p{9.4cm}l@{}}
\toprule
\textbf{Symbol} & \textbf{Meaning} & \textbf{Where} \\
\midrule
\multicolumn{3}{@{}l}{\textit{Confidence and uncertainty}} \\
$\upre$    & Pre-routing confidence: a cheap readiness score read before any tier runs. & \Cref{ch:upre} \\
$\utoken$  & Token-level confidence: the geometric mean of the per-token probabilities of the generated answer. & \Cref{ch:verifier} \\
$\udrop$   & Dropout disagreement: instability of the answer across resampled forward passes. & \Cref{ch:verifier} \\
$\uint$    & Internal confidence: $0.5\,\utoken + 0.5\,(1-\udrop)$, the model's self-assessed certainty. & \Cref{ch:verifier} \\
$\cconv$   & Encoder convergence ratio: how settled the late-layer representation is, a component of $\upre$. & \Cref{ch:upre} \\
$\savg$    & Self-consistency: semantic agreement across the $M$ resampled reasoning chains. & \Cref{ch:verifier} \\
\addlinespace
\multicolumn{3}{@{}l}{\textit{Grounding and relevance}} \\
$\pentail$ & Reasoning-to-answer entailment: whether the chain logically supports the answer. & \Cref{ch:verifier} \\
$\pgmax$   & Top-passage grounding: entailment of the single best supporting passage (the most forgiving reading). & \Cref{ch:verifier} \\
$\pgmean$  & Mean grounding: average support across the retrieved passages. & \Cref{ch:verifier} \\
$\pgatom$  & Atomic grounding: support of the weakest atomic fact in the answer (the strictest reading). & \Cref{ch:verifier} \\
$\qarel$   & Question-answer relevance: how well the answer addresses the question. & \Cref{ch:verifier} \\
\addlinespace
\multicolumn{3}{@{}l}{\textit{Routing and the retrieval-utility classifier}} \\
$\pik$     & Self-knowledge probe: the model's estimated probability that it already knows the answer. & \Cref{ch:ruc} \\
$\prag$    & Predicted retrieval utility: the classifier's probability that retrieval will help this query. & \Cref{ch:ruc} \\
$\tauruc$  & Retrieval-utility threshold: route to retrieval only when $\prag \ge \tauruc = 0.375$. & \Cref{ch:ruc} \\
\addlinespace
\multicolumn{3}{@{}l}{\textit{Trust score and the four-outcome gate}} \\
$\ustored$ & Calibrated composite trust score: the single honest probability the answer is correct. & \Cref{ch:composite} \\
$\taustore$& Storage threshold: store when $\ustored \ge 0.60$. & \Cref{ch:composite} \\
$\taudefer$& Deferral threshold: defer when $0.45 \le \ustored < 0.60$. & \Cref{ch:composite} \\
$\tautrain$& Training trust cut: only episodes with $\ustored \ge 0.70$ train the model. & \Cref{ch:cycle} \\
$\tauretro$& Retroverification prune floor: drop a stored episode whose re-scored trust falls below $0.50$. & \Cref{ch:memory} \\
$\phisafe$ & Safety floor: send a query to retrieval whenever $\upre < 0.38$. & \Cref{ch:upre} \\
\addlinespace
\multicolumn{3}{@{}l}{\textit{Calibration, learning, and theory}} \\
$\rhomin$  & Retention floor: abort a cycle if any probe ratio falls below $0.93$. & \Cref{ch:cycle} \\
$\alpha$   & Shrinkage weight ($0.6$): the blend of a per-benchmark curve toward the pooled curve; also, in theory, verifier balanced accuracy. & \Cref{ch:composite} \\
$\lambda$  & A decay or penalty rate: recency decay in memory, the blend weight in routing, the anchor strength in training. & \Cref{ch:memory} \\
$M,\,K$    & Self-consistency chains ($M=3$) and dropout passes ($K=5$). & \Cref{ch:verifier} \\
$\pi$      & Pool purity: the fraction of stored episodes that are actually correct. & \Cref{ch:theory} \\
$p$        & Base accuracy: the generator's correctness rate before the gate filters it. & \Cref{ch:theory} \\
$d$        & Cohen's $d$: a standardised effect size, used for verifier discrimination. & \Cref{ch:theory} \\
\bottomrule
\end{tabular}
\caption{Mathematical symbols used in the book.}
\label{tab:glossary-symbols}
\end{table}

\section{Terms}
\label{sec:glossary-terms}

\begin{description}[leftmargin=0pt, labelsep=0.6em, style=nextline]

\item[\caem{}] Confidence-Aware Episodic Memory with Self-Improvement: the three-tier,
verifier-gated, cyclically self-improving architecture this book describes.

\item[Abstain] The honest-refusal outcome of the gate: when an answer is both
low-confidence and ungrounded, the system returns ``I do not know'' rather than guess.

\item[Asymmetric rollback] The safety property that, when a cycle aborts, the model
weights revert but the episodic memory is kept. Training can be undone; remembered
knowledge is not.

\item[Atomic-fact decomposition] Splitting an answer into its component factual claims
so each can be checked for grounding separately; the weakest one drives $\pgatom$.

\item[Calibration] Replacing a raw score with a probability whose value means what it
says, so that a stored score of $0.60$ corresponds to a sixty-percent chance of being
correct.

\item[Calibration fold] A held-out, leakage-disjoint slice of labelled answers on which
the composite and the pre-routing temperature are fitted; separate from the evaluation
fold.

\item[Cold start] Cycle zero: the full architecture running before any self-improvement
fine-tuning.

\item[Composite] The calibrated fusion of the verifier's nine signals into the single
trust score $\ustored$, via per-signal isotonic curves, per-benchmark shrinkage, and a
log-odds boost.

\item[Composite evaluation score (CES)] A headline scalar: the geometric mean of five
axes (accuracy, epistemic quality, retention, calibration, verifier reliability), so a
collapse on any one axis drags the whole score down.

\item[Composite hallucination metric (CHM)] The honesty axis: the equal-weighted mean
of eight measured failure-mode rates, lower being better.

\item[Confident confabulation] The costliest failure mode: a wrong answer the system was
nonetheless confident enough to store.

\item[Consolidation] The cycle-boundary merge of near-duplicate stored episodes onto a
single highest-trust representative, guarded by answer agreement and a trust-spread cap.

\item[Defer] The gate outcome for a promising but not-yet-trusted answer: it is parked in
a buffer for reconsideration by a later, better model.

\item[Discard] The gate outcome for a low-confidence answer that is dropped silently,
distinct from an abstention because some grounding exists.

\item[Early-exit] The confabulation gate inside the verifier: when internal confidence is
high but grounding is near zero, the answer is discarded and regenerated.

\item[Episodic memory] The growing store of verified question-answer episodes, embedded
for similarity search and served by the cheap memory tier.

\item[Exact match (EM)] The headline accuracy axis: a normalised answer either matches a
gold answer or does not.

\item[Four-outcome gate] The decision step that reads $\ustored$ and commits to one of
store, defer, abstain, or discard.

\item[Isotonic regression] A monotone, piecewise-constant fit that turns a raw signal
into a calibrated probability of correctness, learning its direction from the data.

\item[Low-rank adaptation (LoRA)] The fine-tuning method: the backbone is frozen and only
a small low-rank adapter trains, which bounds catastrophic forgetting.

\item[Pool reweighting] The temperature-mixed rebalancing of the training pool each cycle
so that no benchmark or answer type dominates the fine-tune.

\item[Purity theorem] The result that a verifier only modestly better than chance lifts
the stored pool's purity above the generator's base accuracy, so an imperfect verifier
still yields a clean memory.

\item[Retention guard] The check that runs three held-out probes after each fine-tune and
aborts the cycle if the worst probe falls below $\rhomin$.

\item[Retrieval-utility classifier (RUC)] The learned gate that predicts whether retrieval
will help a query and routes accordingly, withdrawing grounding where it amplifies the
wrong answer.

\item[Retroverification] The cycle-boundary re-scoring of every stored episode under the
updated model, pruning any whose trust has fallen below $\tauretro$.

\item[Router] The dispatcher that sends each query to the cheap memory tier, the zero-shot
tier, or the retrieval-augmented tier, using similarity, stored trust, and $\upre$.

\item[Self-improvement loop] The outer cycle that fine-tunes the model on its own verified
answers, re-checks retention, and re-scores memory at each boundary.

\item[Store] The gate outcome for a trusted answer: it is admitted to episodic memory as a
verified episode.

\item[Tier 1 / Tier 2 / Tier 3] The three execution paths: a precision-first memory-direct
cache, zero-shot generation, and retrieval-augmented generation.

\item[Training versus transfer] The panel split: three benchmarks the loop trains on and
two it never trains on, the latter measuring generalisation.

\item[Two clocks] The architecture's two timescales: a fast per-query pipeline that only
reads state, and a slow per-cycle loop that is the only thing allowed to change it.

\end{description}

This closes the book. The architecture has been motivated, mapped, built stage by
stage, proven where proof was available, measured across a real six-cycle trajectory,
priced component by component, and reduced to a registry of constants and a glossary of
terms. Every name and number in the preceding chapters can now be traced to its
definition here or to the live code and report the book was built from.
